\documentclass[12pt]{book}
\usepackage[utf8]{inputenc}
\usepackage[T1]{fontenc}
\usepackage{lmodern}
\usepackage{hyperref}
\usepackage{tipa}          % IPA symbols
\usepackage{booktabs}       % Professional tables
\usepackage{amsmath}        % Math for energy formulas
\usepackage{amssymb}
\usepackage{geometry}
\geometry{margin=1in}

\title{Grammar of Lacyo\\[1em]\large A Computationally Optimized Fusional Language\\Derived from the Romance Family}
\author{The Lacyo Project}
\date{Specification v2.0 --- March 2026}

\begin{document}
\maketitle
\tableofcontents

%% ====================================================================
\chapter{Scope and Design Philosophy}
%% ====================================================================

\section{What Lacyo Is}

Lacyo is a constructed language whose phonology, lexicon, and fusional morphology
are derived by computational optimization over the Romance language family.
It is not designed by intuition; it is \emph{calculated}.

The core principle: treat language design as a constrained optimization problem.
A simulated annealing engine searches the combined lexical and morphological space
of 18 Romance languages to find the configuration that minimizes total linguistic
\emph{energy}---a weighted sum of syllable count, phoneme inventory size,
morphological collision, and phonotactic violation.

The human role is to define \emph{constraints} (what is legal) and
\emph{weights} (what is preferred). The machine finds the \emph{forms}.

\section{Design Axioms}

The following axioms are ordered by priority. When they conflict, higher-priority
axioms prevail.

\begin{enumerate}
  \item \textbf{Minimize syllable count.}
        Roots should be as short as possible, measured in phonemic syllables
        (IPA), not orthographic characters. This axiom outranks every other
        preference, including source-language balance and phoneme inventory.
  \item \textbf{One morpheme per meaning-family.}
        Do not keep multiple allomorphs for the same word when a shorter or
        equal-length \emph{attested} form exists (e.g.\ \texttt{gato}/\texttt{gatto}
        collapse to one \texttt{gato}). Never invent a stem by clipping a
        theme vowel: \texttt{gat} is Catalan, not a legal reduction of Spanish
        \texttt{gato}. Ties in syllable count are broken by how many source
        forms the surviving candidate covers.
  \item \textbf{Zero grammatical collision.}
        No two distinct morphological forms within the same paradigm may be
        homophonous. The fusional endings must be fully distinctive.
  \item \textbf{Phonotactic legality.}
        Every morpheme, root, and derived form must conform to the syllable
        structure rules defined in Chapter~\ref{ch:phonology}.
  \item \textbf{Romance recognizability.}
        Where axioms 1--4 permit multiple solutions, prefer forms that are
        recognizable to speakers of at least one Romance language.
        This is a soft preference, not a hard constraint.
  \item \textbf{Fusional morphology.}
        Grammar is encoded in endings that simultaneously express multiple
        categories (e.g., person + number + tense), not by stacking
        one-feature-per-morpheme (agglutinative) or by word order alone
        (isolating).
\end{enumerate}

\section{What Lacyo Is Not}

\begin{itemize}
  \item It is not a pidgin or simplified Romance. It has its own phonology
        and morphological system.
  \item It is not a relexification. Roots are selected, then \emph{adapted}
        to conform to Lacyo phonotactics. ``eau'' is not a Lacyo word;
        its Lacyo form is the IPA-minimal adaptation that fits the phonotactic
        template.
  \item It is not manually designed. If a form cannot be justified by the
        energy function, it does not belong in the language.
\end{itemize}


%% ====================================================================
\chapter{Phonology}
\label{ch:phonology}
%% ====================================================================

\section{Phoneme Inventory}

The phoneme inventory is the output of a cross-Romance frequency analysis.
Phonemes are selected by the optimizer to maximize lexical coverage while
minimizing total inventory size.

The inventory below is the \textbf{initial candidate set}---the search space
from which the optimizer selects. The final inventory is a subset determined
by which roots survive annealing.

\subsection{Consonants (18 candidates)}

\begin{table}[h]
\centering
\caption{Consonant Inventory}
\label{tab:consonants}
\begin{tabular}{@{}lcccccc@{}}
\toprule
             & Bilabial & Labiodental & Alveolar & Postalveolar & Palatal & Velar \\
\midrule
Plosive      & p \quad b &             & t \quad d &              &         & k \quad g \\
Nasal        & m         &             & n         &              &         &           \\
Fricative    &           & f \quad v   & s \quad z & \textipa{S}  &         &           \\
Affricate    &           &             &           & \textipa{tS} &         &           \\
Lateral      &           &             & l         &              &         &           \\
Tap/Trill    &           &             & r         &              &         &           \\
Approximant  &           &             &           &              & j       & w         \\
\bottomrule
\end{tabular}
\end{table}

\noindent \textbf{Total}: 18 consonant phonemes.\\
\textbf{Design notes}:
\begin{itemize}
  \item Voicing distinction preserved for stops (p/b, t/d, k/g) and fricatives
        (f/v, s/z) to maintain Romance recognizability.
  \item \textipa{/S/} (\texttt{sh}) retained: present in French, Portuguese,
        Italian dialectal, and high frequency across Romance loans.
  \item \textipa{/tS/} (\texttt{ch}) retained: present in Italian, Spanish,
        Portuguese, Romanian.
  \item Single rhotic /r/ (tap or trill, free variation). No uvular \textipa{/K/}.
        This is a deliberate aesthetic choice: the ``futuristic'' register
        favors alveolar articulation.
  \item Glides /j/ and /w/ function as consonants in onset position,
        as vowels in nucleus (see phonotactics).
  \item \textbf{Excluded}: \textipa{/K/} (uvular, too marked),
        \textipa{/J/} (palatal nasal---merged to /nj/ sequence),
        \textipa{/L/} (palatal lateral---merged to /lj/),
        /h/ (too weak to be distinctive in a minimal system),
        \textipa{/T/}, \textipa{/D/} (absent from most Romance).
\end{itemize}

\subsection{Vowels (5 phonemes)}

\begin{table}[h]
\centering
\caption{Vowel Inventory}
\label{tab:vowels}
\begin{tabular}{@{}lccc@{}}
\toprule
       & Front & Central & Back \\
\midrule
Close  & i     &         & u    \\
Mid    & e     &         & o    \\
Open   &       & a       &      \\
\bottomrule
\end{tabular}
\end{table}

\noindent \textbf{Total}: 5 vowel phonemes.\\
\textbf{Design notes}:
\begin{itemize}
  \item Classic 5-vowel system shared by Spanish, Italian, Latin, Japanese,
        Swahili, and many other ``clean-sounding'' languages.
  \item No nasal vowels (French \~{a}, \~{o} are denasalized).
  \item No front rounded vowels (French /y/, /\o{}/ are unrounded).
  \item No vowel length distinction. No tonal distinction.
  \item Mid vowels /e/ and /o/ are always tense (never lax). No
        \textipa{/E/}--\textipa{/e/} or \textipa{/O/}--\textipa{/o/} split.
\end{itemize}

\subsection{Total Inventory Size}

\begin{equation}
|\text{Phonemes}| = 18_C + 5_V = 23
\end{equation}

This is the \emph{candidate ceiling}. The optimizer's job is to find lexica
that use as few of these 23 as possible. A lexicon using only 15 consonants
and 5 vowels scores better than one using all 18+5.

\subsection{Diphthongs}

Diphthongs are \textbf{unit phonemes}, not sequences of two vowels.
The following rising and falling diphthongs are attested in the Romance
source data and available in the Lacyo search space:

\begin{itemize}
  \item \textbf{Rising} (glide + vowel): /ja/, /je/, /jo/, /ju/,
        /wa/, /we/, /wi/, /wo/
  \item \textbf{Falling} (vowel + glide): /aj/, /ej/, /oj/,
        /aw/, /ew/, /ow/
\end{itemize}

\noindent Because diphthongs are unit phonemes, each occupies a single vowel
nucleus slot $V$ in the syllable template. A word like /fwe/ (``fire'') is
monosyllabic: $C V$, not $C V V$.

In orthography, diphthongs are spelled as two characters matching their
component phones: \texttt{ya}, \texttt{we}, \texttt{ay}, etc.
The glides /j/ and /w/ are written \texttt{y} and \texttt{w} respectively.


\section{Phonotactics}
\label{sec:phonotactics}

\subsection{Syllable Structure}

The maximal syllable template is:

\begin{equation}
\sigma = (C_1)(C_2) V (C_3)(C_4)
\end{equation}

\noindent where:
\begin{itemize}
  \item $C_1$ (onset 1): any consonant or glide
  \item $C_2$ (onset 2): restricted to /l/, /r/, /j/, /w/ (liquids and glides only)
  \item $V$ (nucleus): any vowel
  \item $C_3$ (coda): any consonant. Coda clusters of two consonants are
        permitted if the first is a sonorant (/m/, /n/, /l/, /r/) and the
        second is an obstruent or /s/ (e.g.\ /ns/, /nd/, /rt/, /lk/).
\end{itemize}

\subsection{Legal Syllable Types}

Ordered by optimizer preference (lower energy first):

\begin{enumerate}
  \item \textbf{CV} --- \texttt{ta, re, su} (preferred: maximally simple)
  \item \textbf{CVC} --- \texttt{tan, res, sul} (common: one coda)
  \item \textbf{V} --- \texttt{a, e, i} (bare nucleus, rare, only word-initial)
  \item \textbf{CCV} --- \texttt{tra, kre, plu} (onset cluster)
  \item \textbf{CCVC} --- \texttt{tran, kres} (onset cluster + coda)
  \item \textbf{VC} --- \texttt{an, es} (rare, only word-initial)
\end{enumerate}

\subsection{Phonotactic Constraints}

\begin{enumerate}
  \item \textbf{Coda clusters.} A syllable may end in one or two consonants.
        Two-consonant codas must follow sonority sequencing: the first consonant
        must be a sonorant (/m/, /n/, /l/, /r/). Three-consonant codas are
        forbidden. This respects naturalistic Romance codas (e.g.\ /film/,
        /dans/, /port/).
  \item \textbf{No onset triples.} Maximum two consonants in onset.
  \item \textbf{Onset cluster rule.} If $C_1 C_2$ onset, then $C_1$ must be
        an obstruent (p, b, t, d, k, g, f, v, s) and $C_2$ must be a
        liquid or glide (l, r, j, w).
  \item \textbf{Any single consonant is a legal coda.} Lacyo respects its
        Romance source material: final stops, fricatives, nasals, and
        liquids are all permitted (e.g.\ /xef/, /dat/, /mem/).
  \item \textbf{No gemination.} No doubled consonants across syllable boundaries.
  \item \textbf{No hiatus.} Adjacent vowels across syllable boundaries must be
        separated by an epenthetic glide or resolved by the optimizer.
\end{enumerate}

\subsection{Stress}

Stress is \textbf{penultimate} (on the second-to-last syllable) by default.
No exceptions, no stress marking needed in orthography. This is the most
common stress pattern across Romance and requires zero additional cognitive
load.

\subsection{Orthography}

Lacyo uses a \textbf{fully phonemic} Latin-script orthography: one letter
per phoneme, one phoneme per letter. No digraphs, no silent letters,
no ambiguity.

\begin{table}[h]
\centering
\caption{Orthographic Mapping}
\label{tab:orthography}
\begin{tabular}{@{}ll|ll|ll@{}}
\toprule
IPA & Letter & IPA & Letter & IPA & Letter \\
\midrule
/p/ & p & /f/ & f & /n/ & n \\
/b/ & b & /v/ & v & /l/ & l \\
/t/ & t & /s/ & s & /r/ & r \\
/d/ & d & /z/ & z & /j/ & y \\
/k/ & k & /\textipa{S}/ & x & /w/ & w \\
/g/ & g & /\textipa{tS}/ & c & & \\
/m/ & m & & & & \\
\midrule
/a/ & a & /i/ & i & /u/ & u \\
/e/ & e & /o/ & o & & \\
\bottomrule
\end{tabular}
\end{table}

\noindent \textbf{Key decisions}:
\begin{itemize}
  \item \texttt{x} for /\textipa{S}/: avoids the digraph ``sh'', keeps one-to-one mapping.
  \item \texttt{c} for /\textipa{tS}/: repurposed from its ambiguous Romance usage
        (where it represents /k/, /s/, or /\textipa{tS}/ depending on language/context).
        In Lacyo, \texttt{c} always and only means /\textipa{tS}/.
  \item \texttt{k} replaces all uses of ``c'' for /k/ and ``qu'' for /kw/.
        ``qu'' does not exist; write \texttt{kw}.
  \item No \texttt{h}, \texttt{q}, \texttt{j} (as consonant---/j/ is written \texttt{y}).
\end{itemize}


%% ====================================================================
\chapter{Morphology}
\label{ch:morphology}
%% ====================================================================

Lacyo is a \textbf{fusional} language: grammatical information is encoded
in endings that simultaneously express multiple categories. This chapter
defines the categories, the paradigm slots, and the constraints the
optimizer must satisfy when generating endings.

\section{Morphological Overview}

Every Lacyo word consists of:

\begin{equation}
\text{Word} = \text{Prefix}^* + \text{Root} + \text{Ending}
\end{equation}

\begin{itemize}
  \item \textbf{Root}: 1--2 syllables. Carries lexical meaning. Selected by
        the optimizer from Romance cognate pool, then adapted to Lacyo
        phonotactics.
  \item \textbf{Ending}: 1 syllable. Fusional: encodes grammatical categories
        simultaneously. Generated by the optimizer subject to distinctiveness
        and phonotactic constraints.
  \item \textbf{Prefix}: 0--2 derivational prefixes, each 1 syllable.
        Optional. From a closed set of $\leq 15$ prefixes.
\end{itemize}

\section{Noun Morphology}
\label{sec:nouns}

\subsection{Categories Encoded in Noun Endings}

\begin{itemize}
  \item \textbf{Gender}: masculine (M), feminine (F)
  \item \textbf{Number}: singular (SG), plural (PL)
\end{itemize}

\noindent This yields $2 \times 2 = 4$ paradigm slots per noun.
Case is not marked (modern Ibero-Romance). Declension is gender $\times$ number.

\subsection{Declension classes are whole-language templates}

Endings are not invented slot-by-slot. The optimizer picks \emph{one}
Romance declension table (Spanish-like \texttt{-o/-a/-os/-as}, Italian-like
\texttt{-o/-a/-i/-e}, \ldots) and applies it as a block. Mixing \texttt{-os}
from Spanish with \texttt{-i} from Italian inside one class is illegal.

\subsection{Paradigm Template}

\begin{table}[h]
\centering
\caption{Noun Paradigm Slots}
\label{tab:noun-paradigm}
\begin{tabular}{@{}lcc@{}}
\toprule
         & Singular  & Plural \\
\midrule
M        & $e_1$     & $e_2$  \\
F        & $e_3$     & $e_4$  \\
\bottomrule
\end{tabular}
\end{table}

\subsection{Gendered stems may mix sources}

For animates with distinct masculine and feminine citation forms, the two
\emph{stems} are chosen independently. Syllables still win, so Catalan
\texttt{gat} (M, 1$\sigma$) may pair with French \texttt{chatte} (F, 1$\sigma$)
even though they come from different languages. That is allowed for nouns.

Inanimates have one stem and take gender from the declension table only.

The definite article is closed-class Portuguese \texttt{o}/\texttt{a}/\texttt{os}/\texttt{as}
(one vowel, gender-marked). It is \emph{not} chosen from Swadesh: \texttt{el}/\texttt{la}
are longer and redundant next to the \texttt{-o/-a} declension.

\section{Verb Morphology --- no mixed-source paradigms}
\label{sec:verbs}

\subsection{Categories Encoded in Verb Endings}

\begin{itemize}
  \item \textbf{Person}: 1st, 2nd, 3rd
  \item \textbf{Number}: singular, plural
  \item \textbf{Tense}: present (PRS), past (PST), future (FUT)
  \item \textbf{Mood}: indicative (IND), subjunctive (SUBJ)
\end{itemize}

\noindent This yields $3 \times 2 \times 3 \times 2 = 36$ paradigm slots per verb.

\subsection{Conjugation Classes}

A verb has \textbf{one stem} and \textbf{one conjugation template}.
Person/number/tense endings are a block from a single Romance source
(Spanish-like, French-like, \ldots). Do \emph{not} take 1sg from Catalan
and 2sg from French: that kind of mix is legal for noun gender, not for
verbs --- it makes the paradigm unlearnable.

The optimizer may pick among 1--3 conjugation classes, each a whole template.

\subsection{Paradigm Template (Indicative)}

\begin{table}[h]
\centering
\caption{Verb Paradigm --- Indicative Mood}
\label{tab:verb-ind}
\begin{tabular}{@{}lcccccc@{}}
\toprule
         & 1SG & 2SG & 3SG & 1PL & 2PL & 3PL \\
\midrule
PRS.IND  & $v_1$ & $v_2$ & $v_3$ & $v_4$ & $v_5$ & $v_6$ \\
PST.IND  & $v_7$ & $v_8$ & $v_9$ & $v_{10}$ & $v_{11}$ & $v_{12}$ \\
FUT.IND  & $v_{13}$ & $v_{14}$ & $v_{15}$ & $v_{16}$ & $v_{17}$ & $v_{18}$ \\
\bottomrule
\end{tabular}
\end{table}

\begin{table}[h]
\centering
\caption{Verb Paradigm --- Subjunctive Mood}
\label{tab:verb-subj}
\begin{tabular}{@{}lcccccc@{}}
\toprule
         & 1SG & 2SG & 3SG & 1PL & 2PL & 3PL \\
\midrule
PRS.SUBJ  & $v_{19}$ & $v_{20}$ & $v_{21}$ & $v_{22}$ & $v_{23}$ & $v_{24}$ \\
PST.SUBJ  & $v_{25}$ & $v_{26}$ & $v_{27}$ & $v_{28}$ & $v_{29}$ & $v_{30}$ \\
FUT.SUBJ  & $v_{31}$ & $v_{32}$ & $v_{33}$ & $v_{34}$ & $v_{35}$ & $v_{36}$ \\
\bottomrule
\end{tabular}
\end{table}

\noindent All 36 $v_i$ must be distinct within a conjugation class.
Each is 1 syllable, phonotactically legal.

\subsection{Infinitive and Non-Finite Forms}

\begin{itemize}
  \item \textbf{Infinitive}: Root + infinitive marker (1 syllable).
        The infinitive marker is class-identifying (like Latin -\={a}re,
        -\={e}re, -\u{e}re, -\={\i}re).
  \item \textbf{Participle}: Root + participle ending (1 syllable).
        One active, one passive. Total: 2 additional endings.
  \item \textbf{Imperative}: 2SG and 2PL only. Total: 2 additional endings.
\end{itemize}

\noindent Grand total per conjugation class: $36 + 1 + 2 + 2 = 41$ endings.


\section{Adjective Morphology}

Adjectives agree with nouns in \textbf{number only} (no case agreement,
no gender agreement).

\begin{itemize}
  \item Singular form: Root + $a_1$
  \item Plural form: Root + $a_2$
\end{itemize}

\noindent Total: 2 endings. Both must be distinct from each other and from
the noun endings to avoid cross-paradigm collision.


\section{Derivational Morphology}
\label{sec:derivation}

Derivational prefixes and suffixes form new words from existing roots.
Each affix is 1 syllable, phonotactically legal.

\subsection{Derivational Prefixes (Candidate Set)}

The optimizer selects from the following semantic slots. Each slot gets
one 1-syllable form:

\begin{table}[h]
\centering
\caption{Derivational Prefix Candidates}
\label{tab:prefixes}
\begin{tabular}{@{}lll@{}}
\toprule
Slot   & Meaning              & Latin Origin \\
\midrule
RE-    & again, back          & \emph{re-}   \\
KON-   & with, together       & \emph{cum-}  \\
DE-    & down, away, reversal & \emph{d\={e}-} \\
EKS-   & out of               & \emph{ex-}   \\
IN-    & in, into             & \emph{in-}   \\
TRA-   & across, through      & \emph{tr\={a}ns-} \\
PRE-   & before               & \emph{prae-} \\
POS-   & after                & \emph{post-} \\
SU-    & under                & \emph{sub-}  \\
EN-    & between              & \emph{inter-} \\
\bottomrule
\end{tabular}
\end{table}

\noindent The forms shown (RE-, KON-, etc.) are human-proposed defaults.
The optimizer may modify these to minimize phoneme inventory, subject to
mutual distinctiveness and phonotactic legality.

\subsection{Derivational Suffixes (Candidate Set)}

\begin{table}[h]
\centering
\caption{Derivational Suffix Candidates}
\label{tab:suffixes}
\begin{tabular}{@{}lll@{}}
\toprule
Slot   & Function                     & Latin Origin \\
\midrule
-ON    & action / process (nominalizer) & \emph{-ti\={o}} \\
-AL    & relating to (adjectivizer)     & \emph{-\={a}lis} \\
-AN    & agent (doer)                   & \emph{-ant-/-ent-} \\
-UR    & result / product               & \emph{-\={u}ra} \\
-EN    & capable of / -able             & \emph{-bilis} \\
-OS    & full of / -ous                 & \emph{-\={o}sus} \\
-AD    & state / quality                & \emph{-it\={a}s} \\
-IS    & abstract noun                  & \emph{-sis/-tis} \\
-IV    & tending to (adjectivizer)      & \emph{-\={\i}vus} \\
-ER    & instrument / tool              & \emph{-\={a}rium} \\
\bottomrule
\end{tabular}
\end{table}

\noindent Derivation stacking: a word may have at most 1 prefix + root + 1 derivational suffix + 1 inflectional ending:

\begin{equation}
\text{MaxWord} = \text{Prefix} + \text{Root} + \text{DerivSuffix} + \text{Ending}
= 1\sigma + 2\sigma + 1\sigma + 1\sigma = 5 \text{ syllables max}
\end{equation}

Most words will be 2--3 syllables (root + ending).


%% ====================================================================
\chapter{The Energy Function}
\label{ch:energy}
%% ====================================================================

This chapter defines the formal optimization objective. The simulated
annealing engine minimizes this function.

\section{Total Energy}

\begin{equation}
E_{\text{total}} = E_{\text{root}} + E_{\text{norm}} + E_{\text{end}} + E_{\text{coll}} + E_{\text{tact}} + E_{\text{dist}}
\end{equation}

Phoneme-inventory size is \textbf{not} an objective ($w_{\text{phon}} = 0$).
The 23-phoneme set is a ceiling, not a budget.

\subsection{Root Syllable Cost}

\begin{equation}
E_{\text{root}} = w_{\text{syl}} \cdot \sum_{r \in \text{Roots}} \text{syl}(r)
\end{equation}

where $\text{syl}(r)$ counts phonemic syllables (IPA-based, not orthographic).
$w_{\text{syl}} = 1000$.

\subsection{Morpheme Normalization}

\begin{equation}
E_{\text{norm}} = w_{\text{norm}} \cdot \sum_{r \in \text{Roots}} (N_{\text{src}} - \mathrm{support}(r))
\end{equation}

$\mathrm{support}(r)$ is how many source-language forms of that concept have
$r$ as a prefix or exact match after Lacyo adaptation. $N_{\text{src}} = 5$
(fr, es, it, pt, ca). $w_{\text{norm}} = 10$ --- enough to break syllable ties
(Catalan \texttt{gat} and French \texttt{xa} are both 1$\sigma$ attested forms)
and far below one extra syllable ($w_{\text{syl}} = 1000$).

Before search, candidates that are the same adapted string or edit-distance
$\leq 1$ are collapsed to a single \emph{attested} representative (shortest,
then highest support). Prefix clipping is not used: a 2$\sigma$ \texttt{gato}
does not spawn a 1$\sigma$ \texttt{gat}.

\subsection{Ending Cost}

\begin{equation}
E_{\text{end}} = w_{\text{end}} \cdot \sum_{e \in \text{Endings}} \text{syl}(e)
\end{equation}

Penalizes multi-syllable endings (they should all be 1 syllable, but the
optimizer must still select the shortest legal forms). $w_{\text{end}} = 200$.

\subsection{Collision Penalty}

\begin{equation}
E_{\text{coll}} = w_{\text{coll}} \cdot |\{(i,j) : e_i = e_j, i \neq j, \text{same paradigm}\}|
\end{equation}

Counts pairs of homophonous endings within the same paradigm (noun declension
or verb conjugation). $w_{\text{coll}} = 100{,}000$ (hard penalty).

\subsection{Phonotactic Violation Penalty}

\begin{equation}
E_{\text{tact}} = w_{\text{tact}} \cdot \sum_{m \in \text{Morphemes}} \text{violations}(m)
\end{equation}

where $\text{violations}(m)$ counts phonotactic rule violations in morpheme $m$.
$w_{\text{tact}} = 500{,}000$ (effectively infinite---violations are illegal).

\subsection{Ending Distinctiveness Penalty}

\begin{equation}
E_{\text{dist}} = w_{\text{dist}} \cdot \sum_{\substack{(i,j) \\ \text{same paradigm}}} \max(0, \; \theta - d(e_i, e_j))
\end{equation}

where $d(e_i, e_j)$ is the phonemic edit distance between endings $e_i$ and $e_j$,
and $\theta = 2$ is the minimum acceptable distance. Endings that differ by
only one feature are penalized. $w_{\text{dist}} = 500$.

This prevents ``confusable'' endings (e.g., \textipa{/-as/} and \textipa{/-at/})
that technically avoid collision but would be error-prone in natural use.


\section{Energy Weights Summary}

\begin{table}[h]
\centering
\caption{Energy Function Weights}
\label{tab:weights}
\begin{tabular}{@{}llrl@{}}
\toprule
Term              & Symbol             & Weight & Rationale \\
\midrule
Root syllables    & $w_{\text{syl}}$   & 1,000  & Primary objective \\
Morpheme support  & $w_{\text{norm}}$  & 10     & Tie-break: one stem, not four \\
Phoneme inventory & $w_{\text{phon}}$  & 0      & Not an objective \\
Ending syllables  & $w_{\text{end}}$   & 200    & Keep endings short \\
Collision         & $w_{\text{coll}}$  & 100,000 & Hard: no homophony \\
Phonotactic viol. & $w_{\text{tact}}$  & 500,000 & Hard: must be legal \\
Distinctiveness   & $w_{\text{dist}}$  & 500    & Soft: avoid confusion \\
\bottomrule
\end{tabular}
\end{table}


%% ====================================================================
\chapter{Lexical Development Pipeline}
\label{ch:lexicon}
%% ====================================================================

\section{Source Data}

Lexical candidates are extracted from Wiktionary XML dumps across 18 Romance
languages (see project README for the full language list and dump sizes).

For each \emph{concept} (semantic slot), candidates are the cognate translations
from each source language, along with their IPA transcriptions.

\section{Pipeline Stages}

\begin{enumerate}
  \item \textbf{Extract}: Parse Wiktionary dumps, extract word + IPA pairs,
        group by concept (via translation links).
  \item \textbf{Phonemicize}: Convert all candidates to Lacyo phoneme space
        using grapheme-to-phoneme (G2P) mapping. Phonemes outside the Lacyo
        inventory are mapped to nearest equivalents (e.g., French
        \textipa{/K/} $\to$ /r/, French /y/ $\to$ /u/).
  \item \textbf{Adapt}: Apply phonotactic constraints. If a candidate violates
        syllable rules, attempt repair (epenthesis, cluster simplification).
        If unrepairable, discard.
  \item \textbf{Optimize}: Feed adapted candidates into the SA engine.
        The engine selects one candidate per concept, minimizing
        $E_{\text{total}}$.
  \item \textbf{Inflect}: Apply fusional endings (also optimizer-generated)
        to each root.
  \item \textbf{Validate}: Check the complete lexicon for inter-word collisions,
        phonotactic violations, and minimum edit distance between
        semantically unrelated words.
\end{enumerate}


\section{Root Adaptation Rules}

When a Romance candidate does not fit Lacyo phonotactics:

\begin{enumerate}
  \item \textbf{Consonant cluster simplification}: Illegal onset/coda clusters
        are reduced. The consonant with lower cross-Romance frequency is dropped.
        Example: /ks/ coda $\to$ /s/.
  \item \textbf{Vowel epenthesis}: If simplification produces an illegal sequence,
        insert /e/ (the default epenthetic vowel). Example: /str/ onset $\to$
        /ser/ (but /tr/ is legal, so /str/ $\to$ /str/ with /s/ moved to
        prior syllable coda if possible).
  \item \textbf{Phoneme mapping}: Non-Lacyo phonemes map as follows:
        \begin{itemize}
          \item \textipa{/K/} (French uvular) $\to$ /r/
          \item /y/ (French front rounded) $\to$ /u/
          \item /\o{}/ (French mid rounded) $\to$ /o/
          \item Nasal vowels $\to$ V + /n/ (e.g., \textipa{/\~{a}/} $\to$ /an/)
          \item \textipa{/\textglotstop/} $\to$ deleted
          \item \textipa{/T/}, \textipa{/D/} (if any) $\to$ /t/, /d/
        \end{itemize}
\end{enumerate}


%% ====================================================================
\chapter{Syntax}
\label{ch:syntax}
%% ====================================================================

\section{Word Order}

Default word order is \textbf{SVO} (Subject--Verb--Object), the dominant
order in modern Romance languages. Because Lacyo has case marking, word
order can be varied for pragmatic emphasis (topicalization, focus), but
SVO is the unmarked/default.

\section{Noun Phrase Structure}

\begin{equation}
\text{NP} = (\text{Det}) + \text{Noun}_{\text{case,num}} + (\text{Adj}_{\text{num}})^*
\end{equation}

Adjectives follow the noun (Romance default). Determiners (articles,
demonstratives) are uninflected particles (isolating behavior at the
function-word level).

\section{Verb Phrase Structure}

\begin{equation}
\text{VP} = (\text{Adv})^* + \text{Verb}_{\text{person,num,tense,mood}} + (\text{NP}_{\text{acc}})
\end{equation}

Since person and number are encoded in the verb ending, subject pronouns
are optional (pro-drop), as in Spanish and Italian.

\section{Subordination}

Subordinate clauses are introduced by uninflected particles (complementizers).
These particles are monosyllabic, selected by the optimizer from a small
closed set.

\section{Negation}

Pre-verbal negation particle. Single syllable, optimizer-selected.

\section{Questions}

Particle-initial for yes/no questions. Wh-words are monosyllabic roots
with case endings.


%% ====================================================================
\chapter{Computational Architecture}
\label{ch:architecture}
%% ====================================================================

\section{Two-Stage Optimization}

\subsection{Stage 1: Root Selection}

For each concept in the lexicon:
\begin{enumerate}
  \item Gather all Romance cognate candidates with IPA transcriptions.
  \item Map each to Lacyo phoneme space.
  \item Apply phonotactic repair.
  \item Score by phonemic syllable count + phoneme inventory contribution.
  \item Feed into SA engine for global optimization.
\end{enumerate}

\subsection{Stage 2: Ending Generation}

Given the grammatical paradigm slots defined in Chapter~\ref{ch:morphology}:
\begin{enumerate}
  \item Enumerate all phonotactically legal 1-syllable strings from the
        Lacyo phoneme inventory.
  \item Assign endings to paradigm slots.
  \item Score by total phoneme cost + collision + distinctiveness.
  \item SA engine searches for optimal assignment.
\end{enumerate}

Stage 2 is run \emph{after} Stage 1, because the phoneme inventory
available for endings depends on which phonemes the roots have already
committed to using.

\section{Genome Format}

The complete state of a Lacyo lexicon is stored in a binary \texttt{.clatin}
file (msgpack serialization):

\begin{verbatim}
{
  "version": "2.0",
  "phoneme_inventory": ["p", "b", "t", ...],
  "roots": {
    "concept_water": {
      "ipa": ["a", "k", "w", "a"],
      "orthography": "akwa",
      "source_lang": "it",
      "source_word": "acqua",
      "syllables": 2
    },
    ...
  },
  "noun_endings": {
    "class_1": {
      "nom_sg": "a", "nom_pl": "as",
      "acc_sg": "on", "acc_pl": "os",
      "gen_sg": "i", "gen_pl": "is"
    },
    ...
  },
  "verb_endings": { ... },
  "adj_endings": { ... },
  "prefixes": { ... },
  "suffixes": { ... },
  "energy_score": 26210,
  "energy_breakdown": { ... }
}
\end{verbatim}

\section{SA Hyperparameters}

\begin{table}[h]
\centering
\caption{Simulated Annealing Configuration}
\begin{tabular}{@{}lr@{}}
\toprule
Parameter                & Value \\
\midrule
Initial temperature      & 10,000 \\
Cooling rate             & 0.9999 \\
Minimum temperature      & 0.01 \\
Expected iterations      & $\sim$200,000 \\
Swap strategy (Stage 1)  & Random concept, random candidate swap \\
Swap strategy (Stage 2)  & Random paradigm slot, random ending swap \\
\bottomrule
\end{tabular}
\end{table}


%% ====================================================================
\chapter{Example Derivations}
\label{ch:examples}
%% ====================================================================

These examples illustrate how the system produces Lacyo forms from Romance
source material. The specific forms shown are \emph{illustrative}---the
actual output depends on the optimizer's run.

\section{Root Selection Example: ``water''}

\begin{tabular}{@{}llrl@{}}
\toprule
Source & IPA & Syllables & Adapted \\
\midrule
FR: eau       & \textipa{/o/}          & 1 & \texttt{o}     \\
ES: agua      & \textipa{/a.gwa/}      & 2 & \texttt{agwa}  \\
IT: acqua     & \textipa{/ak.kwa/}     & 2 & \texttt{akwa}  \\
PT: \'agua    & \textipa{/a.gwA/}      & 2 & \texttt{agwa}  \\
\bottomrule
\end{tabular}

\medskip\noindent
French ``eau'' adapts to monosyllabic \texttt{o}. This scores best on
syllable count. But it's only 1 phoneme, which may cause collision risk.
The optimizer weighs this against keeping \texttt{akwa} (2 syllables,
but common across 3 sources and phonotactically clean).

\section{Fusional Ending Example}

Suppose the optimizer assigns noun class~1 endings:

\begin{tabular}{@{}lcc@{}}
\toprule
         & Singular  & Plural \\
\midrule
NOM      & \texttt{-a}   & \texttt{-as} \\
ACC      & \texttt{-on}  & \texttt{-os} \\
GEN      & \texttt{-i}   & \texttt{-is} \\
\bottomrule
\end{tabular}

\medskip\noindent
Then ``water'' (root \texttt{akw-}) inflects as:

\begin{tabular}{@{}lll@{}}
\toprule
Form & Meaning & Lacyo \\
\midrule
NOM.SG & ``water'' (subject) & \texttt{akwa} \\
ACC.SG & ``water'' (object) & \texttt{akwon} \\
GEN.PL & ``of waters'' & \texttt{akwis} \\
\bottomrule
\end{tabular}


\section{Derivation Example}

Root: \texttt{mun} (``share'', from Latin \emph{m\={u}nus}).

\begin{tabular}{@{}llll@{}}
\toprule
Derivation & Lacyo & Gloss & Formation \\
\midrule
bare root & \texttt{mun-a} & share (NOM.SG) & root + ending \\
kon- + root & \texttt{konmun-a} & commune, community & prefix + root + ending \\
kon- + root + -on & \texttt{konmun-on-a} & communication & prefix + root + suffix + ending \\
\bottomrule
\end{tabular}


%% ====================================================================
\chapter{Open Questions}
\label{ch:open}
%% ====================================================================

The following decisions are deferred to optimizer results or future
specification updates:

\begin{enumerate}
  \item \textbf{Number of declension/conjugation classes.}
        The optimizer may converge on 1, 2, or 3 classes. More classes
        allow shorter endings but increase learning complexity.
  \item \textbf{Definite/indefinite articles.}
        Should Lacyo have articles? If so, monosyllabic particles.
        Pro: Romance heritage. Con: adds function words, increases
        sentence length.
  \item \textbf{Copula.}
        Is there a ``to be'' verb, or are predicate adjectives/nouns
        juxtaposed (zero copula)?
  \item \textbf{Relative clauses.}
        Relative pronoun (inflected?) or particle-based strategy?
  \item \textbf{Number system.}
        Base-10 with monosyllabic roots for 0--9? Vigesimal? The optimizer
        could select the most phonemically efficient numeral roots.
  \item \textbf{Passive voice.}
        Fusional (add paradigm slots = more endings to optimize) or
        periphrastic (particle + active form)?
\end{enumerate}


\end{document}
